\documentclass[11pt,a4paper]{article}
\usepackage[utf8]{inputenc}\usepackage[T1]{fontenc}\usepackage[french]{babel}
\usepackage{amsmath,amssymb,amsthm,geometry,hyperref}
\geometry{margin=2.4cm}
\setlength{\emergencystretch}{3em}
\newtheorem{proposition}{Proposition}
\title{Découpage orthogonal DG avec correction de la constante}
\author{Note exploratoire SLFE}\date{16 novembre 2026}
\begin{document}\maketitle
\paragraph{Statut.} Cette construction a été implémentée et testée sur GPU.
Les invariants et l'ordre temporel sont confirmés sur les petits systèmes
étudiés. Les essais d'advection montrent cependant une erreur de découpage
trop importante aux grands pas, particulièrement sur les maillages gradués.
Aucune supériorité de coût ni nouveauté bibliographique n'est revendiquée.

\section{Construction}
Dans les coordonnées normalisées par la masse, soit $A^T=-A$ l'opérateur DG
centré, avec $Ac=0$ et $\|c\|=1$. Le vecteur $c$ représente la constante.
Décomposer $A=\sum_j A_j$ en opérateurs antisymétriques locaux. Le prototype
sépare les blocs diagonaux des cellules et les couplages de faces. Un coloriage
du graphe dual partitionne les faces en ensembles de paires disjointes.
Les exponentielles locales $U_j(t)=\exp(tA_j)$ sont orthogonales, mais elles
ne préservent pas individuellement $c$.

Pour $b=Uc$, $\gamma=c^Tb>-1$ et $K=cb^T-bc^T$, poser
\[
 Q(b,c)=I+K+\frac{K^2}{1+\gamma},\qquad F_j(t)=Q(U_j(t)c,c)U_j(t).
\]
La matrice $Q$ est la rotation du plan engendré par $b,c$ qui envoie $b$ sur
$c$, et agit comme l'identité sur son orthogonal. Il s'agit d'une formule
classique de rotation, utilisée ici comme correction de l'invariant.

\begin{proposition}
$F_j(t)$ est orthogonal, fixe $c$, et vérifie
$F_j(-t)=F_j(t)^T$. Son générateur infinitésimal est
$F_j'(0)=P A_j P$, avec $P=I-cc^T$.
\end{proposition}
\begin{proof}
Les deux premières propriétés découlent de $Q^TQ=I$ et $Qb=c$.
La rotation est équivariante sous changement orthogonal de coordonnées,
et $Q(b,c)^T=Q(c,b)$. Ainsi
\[
 U_j(t)^T Q(c,U_j(t)c)U_j(t)=Q(U_j(t)^Tc,c),
\]
ce qui donne la symétrie temporelle.
En différentiant en zéro, $b'(0)=A_jc$ et
$Q'(0)=c(A_jc)^T-(A_jc)c^T$.
Comme $c^TA_jc=0$ et $A_j^T=-A_j$, on obtient
$F_j'(0)=A_j-cc^TA_j-A_jcc^T=P A_j P$.
\end{proof}

Puisque $Ac=0$, la somme des générateurs corrigés vaut
$\sum_j P A_j P=PAP=A$. Une composition symétrique des $F_j$ est donc
consistante d'ordre deux à dimension fixée. La composition de Yoshida
$S_2(at)S_2((1-2a)t)S_2(at)$, $a=(2-2^{1/3})^{-1}$, porte l'ordre à quatre.
L'orthogonalité, la masse $c^Tv$ et les constantes sont conservées par
chaque facteur, y compris pour les durées négatives.

\section{Éviter les antipodes pour toute durée}
Supposons que $A_j$ agisse seulement sur un ensemble de cellules $E_j$.
La masse de la composante constante sur cet ensemble est
$\mu_j=\|c|_{E_j}\|^2$. Comme $U_j$ est l'identité ailleurs,
\[
 c^TU_j(t)c\geq (1-\mu_j)-\mu_j=1-2\mu_j.
\]
Le prototype subdivise chaque ensemble de blocs disjoints pour imposer
$\mu_j<1/2$. Le dénominateur $1+\gamma$ est donc strictement supérieur à un,
quelle que soit la durée. Cela requiert que chaque bloc élémentaire couvre
moins de la moitié du domaine, condition vérifiée sur les maillages testés.
Ce découpage permet une stabilité sans restriction de pas; il ne donne
aucune borne d'erreur uniforme par rapport au plus petit élément.

\section{Application de rang deux et réalisation GPU}
Les matrices locales sont exponentiées par diagonalisation hermitienne de
$iA_j$ sur GPU. Pour $\delta=b-c$, une forme de l'application de $Q$ est
\[
 Qx=x+c\frac{2(\gamma-1)c^Tx+2\gamma\delta^Tx}{1+\gamma}
       -\delta\frac{2c^Tx+\delta^Tx}{1+\gamma}.
\]
Elle demande deux produits scalaires globaux et une mise à jour de rang deux.
Le noyau \texttt{rotation\_dg.cu} applique les matrices aux cellules ou aux
paires disjointes. Le coloriage, les partitions en masse et les décompositions
spectrales s'exécutent également sur GPU. Les matrices des sous-pas sont
préparées pour une durée fixe, puis capturées dans un graphe CUDA.

\section{Résultat expérimental et limite}
Six contrôles denses couvrent P1 à P3 sur deux maillages. Ils comparent la
composition à une exponentielle indépendante du générateur, à sa transposée
et à son inverse temporel. Les rapports d'erreur sous division du pas par
deux approchent quatre et seize. Le contrôle mémoire CUDA ne détecte pas
d'erreur sur le cas instrumenté.

Vingt-quatre essais P4 utilisent $n=8,16,32$, deux graduations et
$16,32,64,128$ pas jusqu'à $t=0,5$. L'énergie reste conservée, mais les
erreurs sont grandes. Sur $n=32$ fortement gradué, elles restent proches de
$0,16$ pour ces quatre nombres de pas. Huit raffinements supplémentaires sur
$n=8$ confirment le retour vers le régime spatial lorsque le pas diminue;
sur le maillage gradué, l'erreur passe de $0,0929$ à 256 pas à $0,0170$ à
2048 pas. Les temps correspondants vont d'environ $0,29$ à $2,32$ secondes.
La propriété de stabilité ne supprime donc pas le coût de résolution des
commutateurs locaux. Cette piste n'est pas retenue comme remplacement rapide
du transport semi-lagrangien.
\end{document}
